\documentclass[12pt,a4paper]{article}

% ---- Encoding & Font ----
\usepackage{iftex}
\ifPDFTeX
  \usepackage[utf8]{inputenc}
  \usepackage[T1]{fontenc}
\else
  \usepackage{fontspec}  % XeTeX/LuaTeX (tectonic): Unicode nativo (·, —, ¿, ª, §)
\fi
\usepackage[spanish]{babel}
\usepackage{csquotes}

% ---- Page layout ----
\usepackage[top=2.5cm, bottom=2.5cm, left=2.5cm, right=2.5cm]{geometry}
\usepackage{setspace}
\onehalfspacing

% ---- Math ----
\usepackage{amsmath, amssymb, amsthm}
\usepackage{mathtools}
\usepackage{physics}

% ---- Graphics ----
\usepackage{graphicx}
\usepackage{tikz}
\usepackage{float}
\usepackage{caption}

% ---- Tables ----
\usepackage{array}
\usepackage{booktabs}
\usepackage{tabularx}
\usepackage{colortbl}

% ---- Colours ----
\usepackage{xcolor}
\definecolor{notebg}{HTML}{FFF3CD}
\definecolor{noteborder}{HTML}{FFC107}
\definecolor{boxred}{HTML}{F8D7DA}
\definecolor{boxredborder}{HTML}{F5C6CB}

% ---- Boxes ----
\usepackage{mdframed}
\usepackage{tcolorbox}
\tcbuselibrary{most}

\newtcolorbox{keybox}[1][]{{
  colback=blue!5,
  colframe=blue!70!black,
  arc=4pt,
  boxrule=1pt,
  fonttitle=\\bfseries,
  title={#1}}
}

\newtcolorbox{warnbox}[1][]{{
  colback=boxred,
  colframe=boxredborder,
  coltitle=black!75,
  arc=4pt,
  boxrule=1pt,
  fonttitle=\\bfseries,
  title={#1}}
}

\newtcolorbox{notebox}[1][]{
  colback=notebg,
  colframe=noteborder,
  coltitle=black!75,
  arc=4pt,
  boxrule=1pt,
  fonttitle=\bfseries,
  title={#1}
}

% ---- Hyperlinks & TOC ----
\usepackage[hidelinks]{hyperref}

% ---- Title ----
\title{\textbf{Clase 14: Propiedades algebraicas de la integral}\\\large{Resumen y desarrollo}}
\author{Transcripto y expandido --- Curso Cálculo 1 (OpenFING)}
\date{}

\begin{document}
\maketitle
\tableofcontents
\newpage

\section{Introducción}
En la clase se retomó la noción de integral de Riemann y se mostró que el conjunto de funciones integrables forma un \textbf{subespacio vectorial} del espacio de todas las funciones acotadas.  Esta observación permite aplicar la lógica del álgebra lineal dentro del cálculo, en especial al estudiar cómo se comporta la integral frente a sumas y a multiplicaciones por escalares.

\section{Propiedad de aditividad}

\subsection{Enunciado de aditividad}
\begin{keybox}[Propiedad (Aditividad)]
Si $f$ y $g$ son funciones integrables en el intervalo $[a,b]$, entonces
\[
\int_{a}^{b} \bigl(f(x)+g(x)\bigr)\,dx = \int_{a}^{b} f(x)\,dx + \int_{a}^{b} g(x)\,dx.
\]
\end{keybox}

\subsection{Lema preliminar}
Se necesitó un lema que compara los ínfimos y supremos de la suma con la suma de los ínfimos y supremos:
\begin{itemize}
  \item Para cualquier subintervalo $[c,d]\subset[a,b]$:
  \[
  \inf_{x\in[c,d]}\bigl(f(x)+g(x)\bigr) \ge \inf_{x\in[c,d]} f(x) + \inf_{x\in[c,d]} g(x),
  \]
  \[
  \sup_{x\in[c,d]}\bigl(f(x)+g(x)\bigr) \le \sup_{x\in[c,d]} f(x) + \sup_{x\in[c,d]} g(x).
  \]
Este hecho es una consecuencia directa de la \textbf{desigualdad triangular} aplicada a las funciones.
\end{itemize}

\subsection{Demostración de aditividad}
\begin{enumerate}
  \item Se consideró una partición $P$ del intervalo $[a,b]$ y se definieron las sumas inferiores y superiores de $f$, $g$ y $f+g$.
  \item Usando el lema anterior se obtuvo,
  \[
  S_{P}^{\,\inf}(f+g) \ge S_{P}^{\,\inf}(f) + S_{P}^{\,\inf}(g),
  \]
  \[
  S_{P}^{\,\sup}(f+g) \le S_{P}^{\,\sup}(f) + S_{P}^{\,\sup}(g).
  \]
  \item Al pasar al \textbf{supremo} de todas las sumas inferiores y al \textbf{ínfimo} de todas las sumas superiores se obtuvo la cadena de desigualdades
  \[
  \int_{a}^{b} f + \int_{a}^{b} g \le \int_{a}^{b} (f+g) \le \int_{a}^{b} f + \int_{a}^{b} g,
  \]
  lo que implica igualdad.
\end{enumerate}
\begin{keybox}[Resultado clave]
\[\boxed{\int_{a}^{b} (f+g)\,dx = \int_{a}^{b} f\,dx + \int_{a}^{b} g\,dx}\]
\end{keybox}

\section{Propiedad de homogeneidad}

\subsection{Caso $\alpha = 0$}
Si $\alpha = 0$ la función $\alpha f$ es la función nula, cuya integral es 0. Por tanto
\[
\int_{a}^{b} 0\,dx = 0 = 0\cdot \int_{a}^{b} f\,dx.
\]

\subsection{Caso $\alpha > 0$}
\begin{itemize}
  \item Los ínfimos y supremos de $\alpha f$ se multiplican por $\alpha$ sin cambiar de sentido.
  \item Por el mismo razonamiento que en la aditividad se obtiene
  \[
  \int_{a}^{b} (\alpha f)\,dx = \alpha\int_{a}^{b} f\,dx.
  \]
\end{itemize}

\subsection{Caso $\alpha < 0$}
Multiplicar por un número negativo invierte el orden de las desigualdades, de modo que los ínfimos se convierten en supremos y viceversa. Se muestra que
\[
\int_{a}^{b} (\alpha f)\,dx = \alpha\int_{a}^{b} f\,dx,
\]
manteniendo la igualdad porque la integral inferior coincide con la superior para funciones integrables.
\begin{keybox}[Resultado clave]
\[\boxed{\int_{a}^{b} (\alpha f)\,dx = \alpha\int_{a}^{b} f\,dx}\]
\end{keybox}

\section{Consecuencias estructurales}
Del hecho de que la integral sea lineal se deduce que el conjunto
\[
\mathcal{I}=\{f\,|\,f\text{ integrable en }[a,b]\}\
\]
es un \textbf{subespacio vectorial} de $\mathcal{F}=\{f\,|\,f\text{ acotada en }[a,b]\}$.  Además, el operador integral
\[
\mathcal{T}:\mathcal{I}\to\mathbb{R},\qquad \mathcal{T}(f)=\int_{a}^{b} f(x)\,dx
\]
es una \textbf{transformación lineal}.

Esta observación será útil más adelante cuando se estudien operadores como la derivada, que también resultan lineales sobre los subespacios apropiados.

\section{Comentarios finales y vista a la siguiente clase}
\begin{notebox}{Observación}
Se enfatizó que el razonamiento subyacente a las demostraciones es esencial: no basta con citar el resultado, sino que hay que seguir la cadena de desigualdades y justificar cada paso.
\end{notebox}
En la próxima clase se continuará con la \textbf{construcción del Teorema Fundamental del Cálculo}, mostrando cómo la linealidad permite definir la antiderivada y conectar integrales con derivadas.

\end{document}
